\cleardoublepage % memastikan bab baru dimulai di halaman ganjil (kanan)
\chapter{STUDI LITERATUR}
\label{chap:studi-literatur}

% --- II.1 Hipertensi ---
\section{Hipertensi}
Hipertensi secara global didefinisikan sebagai suatu kondisi hemodinamik kronis yang ditandai dengan meningkatnya tekanan darah di arteri secara persisten. Menurut konsensus medis internasional terbaru, hipertensi bukan sekadar angka tekanan darah yang tinggi, melainkan sindrom kardiovaskular progresif yang memicu perubahan struktural dan fungsional pada pembuluh darah serta organ vital \autocite{Reges2020Association}. Dalam konteks fisiologis, kondisi ini terjadi apabila tekanan darah sistolik mencapai sedikitnya 140 mmHg maupun tekanan darah diastolik sedikitnya 90 mmHg, setelah dilakukan pengukuran berulang di klinik \autocite{Reges2020Association}. International Society of Hypertension (ISH) juga mendukung klaim tersebut, menyatakan bahwa hipertensi pada orang dewasa yaitu diatas 18 tahun, umumnya ditegakkan apabila tekanan darah sistolik mencapai $\geq$140 mmHg dan/atau tekanan darah diastolik mencapai $\geq$90 mmHg.\autocite{Unger2020}. Tekanan darah sistolik menunjukkan tekanan yang dihasilkan ketika jantung berkontraksi untuk memompa darah ke seluruh tubuh, sedangkan tekanan darah diastolik mencerminkan tekanan ketika jantung berada pada fase relaksasi di antara dua denyut.\autocite{Whelton2018Guideline}


Penyakit ini sering dikenal sebagai \emph{the silent killer} karena sebagian besar penderita tidak menunjukkan gejala klinis yang spesifik hingga terjadi kerusakan organ target atau \emph{Target Organ Damage} (TOD) seperti hipertrofi ventrikel kiri, gagal ginjal kronis, atau retinopati \autocite{Mills2020Global}. Prevalensi hipertensi terus meningkat secara global, didorong oleh penuaan populasi dan pergeseran gaya hidup, dengan data terbaru menunjukkan bahwa lebih dari 1,28 miliar orang dewasa berusia 30--79 tahun di seluruh dunia hidup dengan hipertensi, dengan dua pertiga di antaranya tinggal di negara berpenghasilan rendah dan menengah \autocite{WHO2023}.

Di Indonesia, hipertensi masih menjadi tantangan kesehatan utama dengan tingkat pengendalian yang rendah. Studi epidemiologi terbaru menunjukkan bahwa prevalensi hipertensi di Indonesia tidak hanya didominasi oleh lansia, tetapi mulai merambah ke populasi usia produktif akibat transisi nutrisi dan penurunan aktivitas fisik \autocite{manongga2024gambaran}. Patofisiologi hipertensi melibatkan interaksi kompleks antara curah jantung dan resistensi vaskular perifer, yang diatur oleh sistem saraf simpatis, sistem renin-angiotensin-aldosteron (RAAS), dan keseimbangan cairan tubuh. Ketidakseimbangan pada salah satu mekanisme ini dapat memicu vasokonstriksi berkepanjangan yang berujung pada elevasi tekanan darah sistemik.


Hipertensi adalah salah satu faktor risiko utama penyakit kardiovaskular dan berkontribusi terhadap tingginya angka morbiditas maupun mortalitas di seluruh dunia \autocite{WHO2023}. Apabila tidak dikendalikan dengan baik, hipertensi dapat menyebabkan berbagai komplikasi serius, seperti penyakit jantung koroner, gagal jantung, stroke, penyakit ginjal kronis, retinopati, serta meningkatkan risiko kematian dini \autocite{WHO2023}. Oleh karena itu, deteksi dini, pengendalian faktor risiko, dan pengelolaan hipertensi secara berkelanjutan menjadi langkah penting untuk mencegah terjadinya komplikasi dan meningkatkan kualitas hidup penderita.
% \subsection{Klasifikasi Hipertensi}
% Klasifikasi tekanan darah merupakan pedoman klinis fundamental yang digunakan untuk menentukan stratifikasi risiko dan strategi intervensi medis. Meskipun terdapat beberapa pedoman internasional, klasifikasi yang paling banyak diadopsi dalam literatur terkini merujuk pada pedoman International Society of Hypertension (ISH) 2020 dan European Society of Hypertension (ESH) 2023 yang berupaya menyederhanakan kategorisasi untuk kemudahan praktik klinis.

% Perbedaan mendasar dalam klasifikasi modern dibandingkan pedoman lama, seperti JNC 7, adalah penekanan pada kategori "Normal-Tinggi" atau pre-hipertensi sebagai sinyal peringatan dini. Klasifikasi ini penting dalam pengembangan model prediksi berbasis kecerdasan buatan karena memungkinkan algoritma untuk mendeteksi pergeseran pola data pasien dari fase normal menuju fase patologis sebelum diagnosis klinis tegak. Berdasarkan panduan European Society of Cardiology (ESC) dan European Society of Hypertension (ESH) pada tahun 2018, klasifikasi hipertensi dibagi menjadi beberapa kategori sebagai berikut:

% \begin{enumerate}[label=\arabic*., leftmargin=*]
%     \item \textbf{Normal}\\
%     Tekanan darah sistolik $<130$ mmHg dan diastolik $<85$ mmHg dianggap sebagai kondisi ideal dengan risiko kejadian kardiovaskular terendah.
    
%     \item \textbf{High-Normal}\\
%     Rentang sistolik 130--139 mmHg dan/atau diastolik 85--89 mmHg. Kategori ini memerlukan pemantauan ketat karena individu dalam kelompok ini memiliki probabilitas tinggi untuk berkembang menjadi hipertensi menetap jika tidak dilakukan modifikasi gaya hidup.
    
%     \item \textbf{Hipertensi Derajat 1}\\
%     Ditandai dengan sistolik 140--159 mmHg dan/atau diastolik 90--99 mmHg. Pada tahap ini, intervensi farmakologis sering kali mulai dipertimbangkan bersamaan dengan perubahan gaya hidup, tergantung pada profil risiko kardiovaskular total pasien.
    
%     \item \textbf{Hipertensi Derajat 2}\\
%     Tekanan sistolik $\ge 160$ mmHg dan/atau diastolik $\ge 100$ mmHg. Kondisi ini memerlukan penanganan segera untuk mencegah kejadian akut seperti stroke atau serangan jantung.
% \end{enumerate}

% Selain klasifikasi berbasis derajat, hipertensi juga diklasifikasikan berdasarkan etiologinya menjadi hipertensi primer (esensial) dan hipertensi sekunder. Hipertensi primer mencakup 90-95\% dari seluruh kasus dan tidak memiliki penyebab tunggal yang dapat diidentifikasi, melainkan hasil dari interaksi multifaktorial genetik dan lingkungan. Sementara itu, hipertensi sekunder disebabkan oleh kondisi medis yang mendasarinya, seperti penyakit parenkim ginjal, gangguan endokrin, atau penggunaan obat-obatan tertentu.

\subsection{Faktor Risiko}
Menurut \textcite{black2014keperawatan}, faktor risiko hipertensi terklasifikasi menjadi dua, yaitu faktor risiko yang tidak dapat dimodifikasi dan faktor risiko yang dapat dimodifikasi. Pembagian serupa juga digunakan dalam konsensus klinis nasional yang mengelompokkan faktor risiko kardiovaskular pada pasien hipertensi \autocite{PERHI2021}. Identifikasi faktor-faktor ini krusial karena interaksi antar variabel tersebutlah yang akan dianalisis oleh model komputasi untuk menghasilkan prediksi yang akurat.

\subsubsection{Faktor Risiko Tidak Dapat Dimodifikasi}
\begin{enumerate}
    \item \textbf{Usia}\\
    Bertambahnya usia merupakan salah satu faktor risiko hipertensi yang tidak dapat dimodifikasi.\autocite{Mills2020Global} Seiring proses penuaan, pembuluh darah mengalami perubahan struktural dan fungsional, seperti berkurangnya elastisitas arteri, peningkatan kekakuan pembuluh darah, serta penurunan kemampuan endotel dalam mempertahankan homeostasis vaskular.\autocite{Benetos2019} Perubahan tersebut menyebabkan resistensi vaskular perifer meningkat sehingga jantung memerlukan tekanan yang lebih besar untuk mengalirkan darah ke seluruh tubuh.\autocite{Benetos2019} Selain itu, berbagai studi epidemiologi menunjukkan bahwa prevalensi hipertensi meningkat secara progresif seiring bertambahnya usia, terutama setelah memasuki usia paruh baya.\autocite{Mills2020Global} Oleh karena itu, individu berusia lanjut umumnya memiliki risiko lebih tinggi mengalami hipertensi dibandingkan kelompok usia yang lebih muda.
    
    \item \textbf{Jenis Kelamin}\\
    Perbedaan jenis kelamin diketahui memengaruhi risiko seseorang mengalami hipertensi \autocite{Reckelhoff2018}. Pada usia dewasa muda, wanita umumnya memiliki tekanan darah yang lebih rendah dibandingkan pria karena hormon estrogen berperan dalam mempertahankan fungsi endotel, meningkatkan vasodilatasi, serta memberikan efek protektif terhadap sistem kardiovaskular \autocite{Reckelhoff2018}. Namun, setelah memasuki masa menopause, kadar estrogen menurun sehingga efek perlindungan tersebut berkurang dan risiko hipertensi pada wanita meningkat \autocite{Ji2020}. Oleh karena itu, jenis kelamin perlu dipertimbangkan sebagai faktor yang berkontribusi terhadap risiko hipertensi.
    
    \item \textbf{Genetik dan Riwayat Keluarga}\\
    Faktor genetik merupakan salah satu faktor risiko hipertensi yang tidak dapat dimodifikasi \autocite{Padmanabhan2015}. Individu yang memiliki riwayat hipertensi pada orang tua atau saudara kandung cenderung memiliki risiko lebih tinggi untuk mengalami hipertensi dibandingkan individu tanpa riwayat keluarga \autocite{Mills2020Global}. Kondisi tersebut berkaitan dengan adanya pewarisan variasi gen yang memengaruhi regulasi tekanan darah, fungsi ginjal, keseimbangan natrium, serta sistem renin--angiotensin--aldosteron \autocite{Padmanabhan2015}. Kombinasi perubahan pada mekanisme tersebut dapat meningkatkan kecenderungan terjadinya peningkatan tekanan darah \autocite{Padmanabhan2015}.
\end{enumerate}

\subsubsection{Faktor Risiko Dapat Dimodifikasi}
\begin{enumerate}
    \item \textbf{Obesitas dan Indeks Massa Tubuh (IMT)}\\
    Terdapat korelasi linier yang kuat antara peningkatan indeks massa tubuh dan tekanan darah, dan obesitas diperkirakan berkontribusi terhadap sebagian besar kasus hipertensi primer \autocite{Parvanova2024}. Jaringan adiposa yang berlebih, terutama lemak viseral dan perirenal, memicu hipertensi melalui aktivasi berlebih sistem saraf simpatis dan sistem renin--angiotensin--aldosteron, kompresi fisik pada ginjal, serta peningkatan reabsorpsi natrium yang berujung pada ekspansi volume cairan \autocite{Parvanova2024}. Organisasi Kesehatan Dunia juga menempatkan kelebihan berat badan dan obesitas sebagai salah satu faktor risiko metabolik utama hipertensi yang dapat dimodifikasi \autocite{WHO2023}.

    \item \textbf{Asupan Garam (Natrium)}\\
    Konsumsi natrium yang berlebihan merupakan salah satu determinan utama peningkatan tekanan darah melalui peningkatan volume cairan ekstraseluler dan disfungsi endotel. Tinjauan sistematis dan meta-analisis terhadap uji acak terkontrol menunjukkan bahwa pengurangan asupan natrium menurunkan tekanan darah secara bergantung dosis, baik pada individu hipertensi maupun normotensi \autocite{Huang2020Sodium}. Oleh karena itu, pembatasan asupan garam harian menjadi salah satu intervensi gaya hidup paling hemat biaya untuk pengendalian hipertensi \autocite{WHO2023}.

    \item \textbf{Kurangnya Aktivitas Fisik}\\
    Gaya hidup sedentari berkaitan dengan penurunan kebugaran kardiorespirasi dan peningkatan resistensi insulin yang meningkatkan risiko hipertensi \autocite{Hayes2022Physical}. Sebaliknya, aktivitas fisik teratur terbukti menurunkan tekanan darah melalui perbaikan fungsi endotel, penurunan resistensi vaskular sistemik, dan penurunan aktivitas simpatis \autocite{Hayes2022Physical}. Aktivitas fisik yang tidak memadai juga termasuk faktor risiko hipertensi yang direkomendasikan untuk dikendalikan oleh Organisasi Kesehatan Dunia \autocite{WHO2023}.

    \item \textbf{Merokok}\\
    Zat kimia dalam rokok, terutama nikotin, memicu pelepasan katekolamin yang menyebabkan vasokonstriksi akut dan peningkatan denyut jantung, sementara paparan kronis menimbulkan inflamasi sistemik tingkat rendah dan kekakuan arteri \autocite{Gao2023Smoking}. Studi berbasis populasi menunjukkan bahwa merokok berkaitan dengan peningkatan risiko hipertensi, dengan risiko yang lebih besar pada perokok berat \autocite{Gao2023Smoking}.

   \item \textbf{Faktor Psikososial dan Stres}\\
    Stres merupakan respons fisiologis tubuh terhadap tuntutan atau ancaman yang, pada kondisi akut, bersifat adaptif melalui aktivasi sumbu \textit{Hypothalamic--Pituitary--Adrenal} (HPA) dan sistem saraf simpatis untuk mempertahankan homeostasis \autocite{russell2019human}. Aktivasi ini memicu pelepasan hormon kortisol serta katekolamin, terutama norepinefrin dan epinefrin, yang secara sementara meningkatkan denyut jantung, curah jantung, dan resistensi vaskular perifer sehingga tekanan darah naik. Pada stres akut, kenaikan tekanan darah tersebut bersifat sementara dan umumnya kembali normal setelah stresor mereda. Namun, ketika paparan stres berlangsung terus-menerus atau kronis, aktivasi kedua sistem tersebut menjadi berkepanjangan dan berpotensi menimbulkan perubahan yang menetap pada regulasi tekanan darah \autocite{Spruill2010}.

    Paparan stres kronis menyebabkan peningkatan kadar kortisol dan katekolamin secara persisten yang memicu vasokonstriksi berkelanjutan, aktivasi sistem \textit{Renin--Angiotensin--Aldosterone System} (RAAS), retensi natrium, serta \textit{vascular remodeling} \autocite{scalco2005hypertension}. Rangkaian perubahan fisiologis ini meningkatkan resistensi vaskular perifer dan volume cairan tubuh sehingga tekanan darah cenderung menetap pada level yang lebih tinggi. Selain melalui jalur fisiologis langsung, stres kronis juga berkontribusi secara tidak langsung dengan mendorong perilaku tidak sehat seperti pola makan buruk, konsumsi alkohol, merokok, dan gangguan tidur yang turut meningkatkan risiko hipertensi \autocite{Spruill2010}. Kombinasi kedua jalur tersebut menjadikan stres psikososial sebagai salah satu faktor risiko hipertensi yang dapat dimodifikasi.

    Bukti epidemiologis mengenai hubungan stres dan hipertensi terus berkembang meskipun belum sepenuhnya konsisten. Tinjauan sistematis oleh \textcite{Sparrenberger2009} menyimpulkan bahwa stres akut hanya menimbulkan kenaikan tekanan darah sementara, tetapi sejumlah studi kohort menunjukkan bahwa stres psikososial kronis serta faktor psikososial lain seperti kecemasan berkaitan dengan peningkatan kejadian hipertensi. Temuan tersebut diperkuat oleh meta-analisis \textcite{Gasperin2009} terhadap studi kohort yang menemukan bahwa individu dengan tingkat stres psikologis tinggi memiliki risiko peningkatan tekanan darah yang lebih besar, meskipun besar efeknya tergolong moderat. Dengan demikian, stres dan faktor psikososial merupakan aspek penting yang perlu dipertimbangkan dalam penilaian risiko hipertensi.
    
    \item \textbf{Diabetes Melitus}\\
    Diabetes melitus dan hipertensi kerap muncul secara bersamaan karena berbagi mekanisme patofisiologis yang sama, meliputi resistensi insulin, disfungsi endotel, dan aktivasi sistem renin--angiotensin--aldosteron \autocite{Petrie2018Diabetes}. Hipertensi dilaporkan sekitar dua kali lebih sering ditemukan pada penderita diabetes dibandingkan individu tanpa diabetes, dan koeksistensi kedua kondisi tersebut secara substansial meningkatkan risiko komplikasi kardiovaskular \autocite{Petrie2018Diabetes}. Oleh karena itu, status diabetes merupakan faktor risiko penting yang perlu diperhitungkan dalam penilaian risiko hipertensi.

    \item \textbf{Kolesterol Tinggi (Dislipidemia)}\\
    Dislipidemia, khususnya kadar \emph{Low-Density Lipoprotein} (LDL) yang tinggi, memicu disfungsi endotel yang menjadi tahap awal aterosklerosis serta gangguan kemampuan pembuluh darah untuk berdilatasi \autocite{Higashi2023Dyslipidemia}. Gangguan fungsi endotel ini dapat berkontribusi terhadap peningkatan tekanan darah, antara lain melalui aktivasi sistem renin--angiotensin--aldosteron, vasokonstriksi, dan retensi natrium \autocite{Higashi2023Dyslipidemia}. Dengan demikian, kadar kolesterol yang tinggi kerap menyertai dan memperberat kondisi hipertensi.

    \item \textbf{Kualitas Tidur}\\
    Kualitas dan durasi tidur yang buruk berkaitan dengan peningkatan risiko hipertensi. Meta-analisis menunjukkan bahwa kualitas tidur subjektif yang buruk berhubungan dengan peningkatan kemungkinan terjadinya hipertensi \autocite{Lo2018Sleep}. Mekanisme yang mendasari mencakup peningkatan tonus simpatis yang menetap serta terganggunya penurunan tekanan darah pada malam hari (\emph{nocturnal dipping}), yang dalam jangka panjang dapat memicu hipertensi \autocite{Lo2018Sleep}. Oleh karena itu, kualitas tidur menjadi faktor gaya hidup yang layak untuk dipertimbangkan dalam penilaian risiko hipertensi.
\end{enumerate}

% --- II.2 Machine Learning ---
\section{Machine Learning}
Machine learning merupakan salah satu bidang dalam kecerdasan buatan yang berfokus pada pengembangan model komputasi agar mampu mengenali pola dari data dan menghasilkan prediksi tanpa bergantung pada aturan yang ditentukan secara eksplisit \autocite{Sarker2021}. Dalam konteks prediksi risiko hipertensi, permasalahan yang dihadapi tergolong klasifikasi biner terawasi (\emph{supervised binary classification}), yaitu memetakan sekumpulan fitur faktor risiko pasien ke dalam salah satu dari dua kelas: berisiko atau tidak berisiko \autocite{Sarker2021}. Data yang digunakan pada penelitian ini bersifat tabular dengan kombinasi fitur numerik dan kategorikal, sehingga pemilihan algoritma difokuskan pada metode yang secara empiris menunjukkan kinerja kuat pada karakteristik data tersebut. Bagian ini membahas algoritma klasifikasi yang digunakan, teknik penanganan ketidakseimbangan kelas, normalisasi fitur, penyetelan hiperparameter, serta metrik yang dipakai untuk mengevaluasi kinerja model.

\subsection{Algoritma Machine Learning}
Bagian ini menguraikan empat algoritma klasifikasi yang dievaluasi dalam penelitian ini, yaitu Logistic Regression sebagai model linear \emph{baseline}, serta tiga model berbasis \emph{ensemble} pohon keputusan: Random Forest, XGBoost, dan CatBoost. Keempatnya dipilih karena secara konsisten dilaporkan berkinerja baik pada masalah klasifikasi data tabular di ranah kesehatan \autocite{Septian2025Prediction}. Logistic Regression diikutkan sebagai pembanding karena sifatnya yang sederhana dan mudah diinterpretasi, sedangkan ketiga model \emph{ensemble} dipertimbangkan karena kemampuannya menangkap hubungan non-linear antar faktor risiko. Perbandingan keempat algoritma bertujuan menemukan model dengan keseimbangan terbaik antara akurasi prediksi dan kestabilan kinerja pada data yang tidak seimbang.

\subsubsection{Logistic Regression}
Logistic Regression adalah metode klasifikasi statistik yang memodelkan probabilitas suatu kejadian biner sebagai fungsi logistik dari kombinasi linear fitur masukan \autocite{Hosmer2013Applied}. Model ini mengestimasi koefisien untuk setiap fitur sehingga besar maupun arah pengaruh masing-masing faktor risiko terhadap peluang hipertensi dapat dibaca secara langsung, misalnya melalui nilai \emph{odds ratio}. Keunggulan utamanya terletak pada kesederhanaan komputasi dan interpretabilitas yang tinggi, sehingga algoritma ini kerap dijadikan \emph{baseline} dalam studi prediksi penyakit \autocite{Zhang2025Machine}. Namun, karena mengasumsikan hubungan linear antara fitur dan \emph{log-odds}, model ini cenderung kurang optimal dalam menangkap interaksi kompleks dan pola non-linear yang umum ditemukan pada data klinis \autocite{Tarimana2024}.

\subsubsection{Random Forest}
Random Forest merupakan algoritma \emph{ensemble} yang menggabungkan banyak pohon keputusan (\emph{decision trees}) untuk menghasilkan prediksi yang lebih akurat dan stabil dibandingkan pohon tunggal \autocite{Breiman2001RandomForests}. Setiap pohon dilatih pada subset data yang dibentuk melalui \emph{bootstrap} (\emph{bagging}) dengan pemilihan fitur acak pada tiap percabangan, sehingga antarpohon memiliki korelasi rendah dan kesalahan individual saling menetralkan. Prediksi akhir diperoleh melalui pemungutan suara mayoritas untuk klasifikasi, yang membuat model relatif tahan terhadap \emph{overfitting} dan \emph{noise}. Mekanisme keacakan ganda ini menjadikan Random Forest andal pada data tabular kesehatan, meskipun kinerjanya pada dataset berskala besar dapat tertinggal dibandingkan metode \emph{boosting} modern \autocite{S2024Impact}.

\subsubsection{Extreme Gradient Boosting (XGBoost)}
Extreme Gradient Boosting (XGBoost) adalah implementasi \emph{gradient boosting} yang membangun pohon keputusan secara berurutan, dengan setiap pohon baru difokuskan untuk memperbaiki kesalahan residual dari gabungan pohon sebelumnya \autocite{Chen2016XGBoost}. Berbeda dengan \emph{gradient boosting} standar, XGBoost menyisipkan regularisasi (\emph{L1} dan \emph{L2}) pada fungsi objektifnya untuk mengendalikan kompleksitas model dan menekan risiko \emph{overfitting} \autocite{Chen2016XGBoost}. Algoritma ini juga dirancang untuk efisiensi komputasi melalui pemrosesan paralel dan penanganan nilai hilang secara internal, sehingga mampu mengolah data dalam jumlah besar dengan cepat. Karakteristik tersebut menjadikan XGBoost salah satu algoritma paling banyak digunakan dan berkinerja tinggi dalam prediksi risiko penyakit berbasis data tabular \autocite{Li2025XGBoost}.

\subsubsection{Categorical Boosting (CatBoost)}
CatBoost adalah algoritma \emph{gradient boosting} yang dikembangkan dengan penekanan khusus pada penanganan fitur kategorikal secara efisien \autocite{Prokhorenkova2018CatBoost}. Untuk menghindari kebocoran informasi target (\emph{target leakage}) yang lazim terjadi pada teknik \emph{target encoding} konvensional, CatBoost menerapkan \emph{ordered target statistics} yang menghitung nilai enkode hanya dari data yang muncul sebelumnya dalam suatu urutan acak. Selain itu, algoritma ini menggunakan \emph{ordered boosting} untuk mengurangi bias pada estimasi gradien sehingga model yang dihasilkan lebih stabil dan tahan terhadap \emph{overfitting}. Keunggulan tersebut membuat CatBoost sering menunjukkan performa kompetitif pada data tabular, termasuk pada studi prediksi penyakit di populasi Indonesia \autocite{Septian2025Prediction}.

\subsection{Class Weighting}
Ketidakseimbangan kelas (\emph{class imbalance}) terjadi ketika jumlah sampel pada suatu kelas jauh lebih sedikit dibandingkan kelas lainnya, kondisi yang umum ditemukan pada data medis karena penderita suatu penyakit biasanya lebih sedikit daripada individu sehat \autocite{he2009}. Ketimpangan ini menyebabkan model cenderung bias terhadap kelas mayoritas dan gagal mengenali kelas minoritas yang justru menjadi fokus utama deteksi dini. \emph{Class weighting} merupakan pendekatan pada tingkat algoritma yang mengatasi masalah tersebut dengan memberikan bobot lebih besar pada kesalahan klasifikasi kelas minoritas di dalam fungsi kerugian (\emph{loss function}). Berbeda dengan metode \emph{resampling} yang mengubah distribusi data, teknik ini mempertahankan data asli sehingga tidak berisiko menghilangkan informasi maupun menghasilkan sampel sintetik yang menyesatkan \autocite{krawczyk2016}.

\subsection{Z-score Normalization}
Z-score normalization adalah teknik penskalaan fitur numerik yang mentransformasikan setiap nilai menjadi skor standar berdasarkan rata-rata dan simpangan baku distribusinya \autocite{Patro2015Normalization}. Hasil transformasi menghasilkan fitur dengan rata-rata nol dan simpangan baku satu, sehingga fitur-fitur dengan rentang nilai yang sangat berbeda, misalnya usia dan indeks massa tubuh, berada pada skala yang setara. Penyeragaman skala ini penting agar tidak ada fitur yang mendominasi proses pelatihan semata-mata karena rentang nilainya lebih besar. Teknik ini khususnya bermanfaat bagi algoritma yang sensitif terhadap skala fitur seperti Logistic Regression, sementara pada model berbasis pohon pengaruhnya cenderung minimal namun tetap menjaga konsistensi tahap praproses \autocite{Patro2015Normalization}.

\subsection{Hyperparameter Tuning}
Hyperparameter adalah parameter konfigurasi model yang nilainya ditetapkan sebelum proses pelatihan dan tidak dipelajari langsung dari data, seperti \emph{learning rate}, jumlah pohon, dan kedalaman maksimum pohon \autocite{Bergstra2012Random}. Hyperparameter tuning merupakan proses pencarian kombinasi nilai hyperparameter yang menghasilkan kinerja model terbaik pada data validasi. Untuk menilai setiap kombinasi secara andal, kinerja umumnya diukur menggunakan validasi silang k-lipatan (\emph{k-fold cross-validation}), yaitu membagi data latih menjadi k bagian dan menggilir salah satunya sebagai data validasi agar estimasi kinerja tidak bergantung pada satu pembagian data tertentu. Pencarian kombinasi tersebut dapat dilakukan dengan beberapa strategi, di antaranya \emph{grid search} yang menelusuri seluruh kombinasi secara sistematis dan \emph{random search} yang mengambil sampel kombinasi secara acak. \textcite{Bergstra2012Random} menunjukkan bahwa \emph{random search} sering kali lebih efisien daripada \emph{grid search} karena mampu menemukan konfigurasi yang baik dengan jumlah percobaan yang lebih sedikit, terutama ketika hanya sebagian hyperparameter yang berpengaruh signifikan terhadap kinerja.

\subsection{Optimasi Threshold}
Model klasifikasi biner umumnya menghasilkan probabilitas atau skor risiko yang perlu diubah menjadi keputusan kelas melalui suatu nilai ambang (\emph{threshold}). Ambang bawaan sebesar 0,5 tidak selalu merupakan titik operasi terbaik, terutama pada data tidak seimbang atau ketika kesalahan pada salah satu kelas lebih merugikan daripada kelas lainnya. Optimasi threshold bertujuan mencari nilai ambang yang paling sesuai dengan prioritas permasalahan, misalnya menyeimbangkan sensitivity dan specificity atau memaksimalkan penangkapan kasus positif. Salah satu metode yang umum digunakan adalah Youden Index ($J$), yang dihitung sebagai selisih antara \emph{true positive rate} dan \emph{false positive rate} ($J = TPR - FPR$); ambang dengan nilai $J$ tertinggi dipilih karena merepresentasikan titik operasi yang paling menjauh dari garis tebakan acak pada kurva ROC \autocite{Habibzadeh2016Cutoff}.

Dalam konteks tugas akhir ini, titik operasi Youden dipilih karena paling sesuai dengan karakteristik permasalahan yang dihadapi, yaitu deteksi dini risiko hipertensi pada data yang tidak seimbang. Secara geometris, nilai $J$ setara dengan jarak vertikal antara suatu titik pada kurva ROC dan garis diagonal tebakan acak, sehingga ambang dengan $J$ tertinggi merupakan titik yang memaksimalkan selisih antara proporsi kasus positif yang berhasil dideteksi (\emph{sensitivity}) dan proporsi kasus negatif yang keliru ditandai sebagai positif (\emph{1 $-$ specificity}) \autocite{Habibzadeh2016Cutoff}. Titik tersebut merepresentasikan keseimbangan terbaik antara kemampuan menangkap individu berisiko dan kemampuan menghindari alarm palsu. Kriteria Youden juga tidak bergantung pada proporsi kelas dalam data sehingga tetap stabil pada kondisi kelas minoritas yang jarang dan tidak tertarik ke arah kelas mayoritas sebagaimana ambang yang dioptimalkan semata-mata berdasarkan \emph{accuracy} \autocite{Perkins2006Youden}. Perlu ditegaskan bahwa kriteria Youden bukan mekanisme utama untuk memprioritaskan kelas positif, melainkan diterapkan di atas model yang sebelumnya telah dibuat sensitif terhadap kelas hipertensi melalui pembobotan kelas (\emph{class weighting}) pada tahap pelatihan. Prioritas untuk menekan kasus terlewat (\emph{false negative}) ditegakkan terutama pada tahap pemodelan, yaitu melalui metrik yang berorientasi pada \emph{recall} serta ambang \emph{recall} minimum yang ditetapkan sejak awal sebagai syarat kelayakan model. Pada sistem skrining hipertensi, kesalahan melewatkan individu berisiko (\emph{false negative}) memang dinilai lebih merugikan karena menunda upaya pencegahan, tetapi alarm palsu (\emph{false positive}) tetap membebani alur rujukan pemeriksaan lanjutan sehingga tidak dapat diabaikan. Oleh karena itu, di antara berbagai titik operasi yang mungkin, ambang Youden dipilih karena memberikan keseimbangan yang reprodusibel dan tidak bergantung pada proporsi kelas; apabila di kemudian hari dikehendaki postur yang lebih agresif dalam menekan \emph{false negative}, ambang keputusan dapat diturunkan untuk memenuhi target \emph{recall} tertentu tanpa melatih ulang model karena keluaran probabilitas tetap dipertahankan.

\subsection{Metrik Evaluasi}
Evaluasi kinerja model klasifikasi memerlukan metrik yang sesuai dengan karakteristik data, terutama pada kasus ketidakseimbangan kelas yang membuat sebagian metrik menjadi kurang representatif. Sebagian besar metrik diturunkan dari \emph{confusion matrix} yang memuat jumlah \emph{true positive} (TP), \emph{true negative} (TN), \emph{false positive} (FP), dan \emph{false negative} (FN). Bagian ini menguraikan metrik-metrik yang digunakan untuk menilai kualitas prediksi model dalam penelitian ini. Penggunaan beberapa metrik secara bersamaan bertujuan memberikan gambaran kinerja yang lebih menyeluruh sekaligus menghindari kesimpulan yang menyesatkan.

Keempat komponen \emph{confusion matrix} tersebut didefinisikan berdasarkan kesesuaian antara kelas hasil prediksi dan kelas sebenarnya, dengan kelas positif pada penelitian ini merepresentasikan individu yang berisiko hipertensi \autocite{Fawcett2006Introduction}. \emph{True positive} (TP) adalah jumlah individu berisiko hipertensi yang diprediksi dengan benar sebagai berisiko, sedangkan \emph{true negative} (TN) adalah jumlah individu yang tidak berisiko dan diprediksi dengan benar sebagai tidak berisiko. Kedua komponen ini merepresentasikan prediksi yang tepat, sehingga semakin besar nilainya menandakan kinerja model yang semakin baik dalam mengenali kedua kelas secara benar.

Sebaliknya, \emph{false positive} (FP) menyatakan individu yang sebenarnya tidak berisiko tetapi diprediksi berisiko, sementara \emph{false negative} (FN) menyatakan individu yang sebenarnya berisiko tetapi diprediksi tidak berisiko \autocite{Sokolova2009Systematic}. Dalam konteks deteksi dini hipertensi, kedua jenis kesalahan ini memiliki konsekuensi yang berbeda: FN berpotensi lebih berbahaya karena individu yang benar-benar berisiko tidak terdeteksi sehingga kehilangan kesempatan untuk memperoleh intervensi dini, sedangkan FP menimbulkan kekhawatiran dan pemeriksaan lanjutan yang sebenarnya tidak diperlukan. Perbedaan konsekuensi ini tidak diatasi dengan membiaskan titik operasi model, melainkan melalui pemilihan metrik yang berorientasi pada \emph{recall} serta penetapan ambang \emph{recall} minimum sebagai syarat kelayakan model pada tahap pemodelan, sehingga kepekaan terhadap kasus positif tetap diutamakan tanpa mengabaikan beban \emph{false positive}. Karena hampir seluruh metrik evaluasi, seperti \emph{recall}, \emph{specificity}, \emph{precision}, dan F1-Score, diturunkan dari keempat komponen tersebut, pemahaman terhadap TP, TN, FP, dan FN menjadi dasar penting dalam menilai kualitas prediksi model \autocite{Sokolova2009Systematic}.

\subsubsection{ROC-AUC}
ROC-AUC (\emph{Area Under the Receiver Operating Characteristic Curve}) mengukur kemampuan model dalam membedakan mana kelas positif dan yang negatif di seluruh nilai ambang (\emph{threshold}) klasifikasi \autocite{Fawcett2006Introduction}. Kurva ROC dibentuk dengan memplot \emph{true positive rate} terhadap \emph{false positive rate}, dan luas area di bawahnya merepresentasikan probabilitas model memberikan skor lebih tinggi pada sampel positif acak dibandingkan sampel negatif acak. Nilai ROC-AUC berkisar dari 0,5 yang setara dengan tebakan acak hingga 1,0 yang menandakan pemisahan sempurna antar kelas. Karena tidak bergantung pada satu ambang tertentu, metrik ini kerap digunakan sebagai ukuran kinerja umum pada klasifikasi biner \autocite{Fawcett2006Introduction}.

\subsubsection{PR-AUC}
PR-AUC (\emph{Area Under the Precision-Recall Curve}) mengukur luas area di bawah kurva yang memplot \emph{precision} terhadap \emph{recall} pada berbagai nilai ambang \autocite{Saito2015Precision}. Metrik ini berfokus pada kinerja model terhadap kelas positif sehingga lebih informatif dibandingkan ROC-AUC ketika kelas positif jauh lebih sedikit daripada kelas negatif. Pada kondisi tersebut, ROC-AUC dapat memberikan kesan optimistis yang menyesatkan, sementara PR-AUC lebih realistis menggambarkan \emph{trade-off} antara ketepatan dan kelengkapan deteksi kasus positif \autocite{Saito2015Precision}. Nilai dasar (\emph{baseline}) PR-AUC bergantung pada proporsi kelas positif dalam data, sehingga interpretasinya harus mempertimbangkan tingkat ketidakseimbangan dataset \autocite{Saito2015Precision}.

\subsubsection{Specificity}
Specificity mengukur proporsi kasus negatif yang berhasil diklasifikasikan dengan benar oleh model, yaitu rasio \emph{true negative} terhadap seluruh kasus yang sebenarnya negatif \autocite{Sokolova2009Systematic}. Metrik ini mencerminkan kemampuan model dalam menghindari \emph{false positive}, yakni menghindari pelabelan individu sehat sebagai berisiko. Dalam konteks skrining hipertensi, specificity yang rendah dapat menimbulkan kekhawatiran berlebih serta pemeriksaan lanjutan yang sebenarnya tidak diperlukan. Oleh karena itu, specificity umumnya dianalisis berdampingan dengan sensitivity untuk menilai keseimbangan kinerja model pada kedua kelas \autocite{Sokolova2009Systematic}.

\subsubsection{Recall (Sensitivity)}
Recall, yang juga dikenal sebagai sensitivity, mengukur proporsi kasus positif yang berhasil dikenali dengan benar oleh model, yaitu rasio \emph{true positive} terhadap seluruh kasus yang sebenarnya positif \autocite{Sokolova2009Systematic}. Metrik ini sangat penting pada aplikasi medis karena kegagalan mendeteksi individu yang benar-benar berisiko (\emph{false negative}) dapat berakibat fatal akibat tertundanya penanganan. Nilai recall yang tinggi memastikan sebagian besar penderita terdeteksi, meskipun konsekuensinya dapat berupa peningkatan jumlah \emph{false positive}. Dalam sistem deteksi dini hipertensi, recall kerap dijadikan salah satu metrik prioritas karena biaya melewatkan kasus positif jauh lebih besar daripada biaya alarm palsu \autocite{Sokolova2009Systematic}.

\subsubsection{F1-Score}
F1-Score adalah rata-rata harmonik (\emph{harmonic mean}) dari precision dan recall yang merangkum keduanya ke dalam satu nilai tunggal \autocite{Sokolova2009Systematic}. Penggunaan rata-rata harmonik membuat metrik ini hanya bernilai tinggi apabila precision dan recall sama-sama tinggi, sehingga model tidak dapat memperoleh skor baik dengan mengorbankan salah satunya. Karakteristik tersebut menjadikan F1-Score lebih representatif dibandingkan accuracy pada data yang tidak seimbang, terutama ketika baik \emph{false positive} maupun \emph{false negative} sama-sama memiliki konsekuensi penting. F1-Score banyak digunakan pada domain medis ketika ketepatan dan kelengkapan deteksi sama-sama diperhitungkan \autocite{Sokolova2009Systematic}.

\subsubsection{Accuracy}
Accuracy mengukur proporsi prediksi yang benar terhadap keseluruhan prediksi yang dibuat oleh model \autocite{Sokolova2009Systematic}. Metrik ini populer karena mudah dipahami dan dihitung, namun dapat menyesatkan pada kasus klasifikasi dengan distribusi kelas yang timpang. Sebagai contoh, pada dataset dengan mayoritas kelas negatif, model yang selalu memprediksi kelas mayoritas tetap memperoleh accuracy tinggi meskipun gagal mengenali satu pun kasus positif. Oleh karena itu, accuracy sebaiknya tidak digunakan sendirian, melainkan dilengkapi metrik lain seperti F1-Score, ROC-AUC, atau PR-AUC agar penilaian kinerja lebih holistik \autocite{Sokolova2009Systematic}.

% --- II.3 Explainable Artificial Intelligence (XAI) ---
\section{Explainable Artificial Intelligence (XAI)}
Penerapan Machine Learning dalam praktik klinis sering terhambat oleh sifat \textit{black box} dari model berkinerja tinggi seperti XGBoost, yang proses pengambilan keputusannya sulit ditelusuri oleh manusia. Explainable Artificial Intelligence (XAI) hadir sebagai solusi untuk memberikan transparansi terhadap keputusan yang dihasilkan oleh sistem cerdas sehingga alasan di balik setiap prediksi dapat dipahami \autocite{Adadi2018}. Transparansi ini krusial di ranah kesehatan karena tenaga medis perlu memahami dasar suatu prediksi sebelum menggunakannya dalam pengambilan keputusan yang menyangkut keselamatan pasien. Di antara berbagai metode XAI yang tersedia, penelitian ini menggunakan Shapley Additive Explanations (SHAP) karena landasan teoretisnya yang kuat serta konsistensi penjelasannya, dibandingkan metode alternatif seperti Local Interpretable Model-agnostic Explanations (LIME) yang penjelasannya cenderung kurang stabil pada eksekusi berulang \autocite{molnar2025interpretable}.

\subsection{Shapley Additive Explanations (SHAP)}
\label{subsec:landasan-shap}
Metode SHAP merupakan pendekatan XAI yang mengadopsi konsep teori permainan kooperatif untuk menghitung kontribusi marginal (\emph{Shapley value}) setiap fitur terhadap hasil prediksi model \autocite{Lundberg2017Unified}. \textcite{Lundberg2017Unified} memperkenalkan SHAP sebagai kerangka terpadu (\emph{unified framework}) yang menyatukan berbagai metode atribusi fitur ke dalam satu kelas metode aditif, dan membuktikan bahwa nilai Shapley merupakan satu-satunya solusi yang secara simultan memenuhi tiga aksioma yang diinginkan, yaitu \emph{local accuracy} (jumlah kontribusi seluruh fitur setara dengan selisih antara prediksi aktual dan prediksi rata-rata model), \emph{missingness}, dan \emph{consistency}. Sifat aditif dan berbasis aksioma inilah yang menjadikan SHAP lebih unggul dibandingkan metode berbasis perturbasi lokal seperti LIME, yang penjelasannya dapat berubah-ubah untuk data yang sama pada setiap iterasi eksekusi \autocite{molnar2025interpretable}. Karena stabilitas dan konsistensi penjelasan merupakan syarat penting dalam sistem pendukung keputusan medis, SHAP dipilih sebagai metode interpretasi utama dalam penelitian ini untuk mengidentifikasi kontribusi tiap faktor risiko terhadap prediksi hipertensi.

Tantangan utama dalam penerapan nilai Shapley terletak pada beban komputasinya. Perhitungan nilai Shapley secara eksak pada kasus umum menuntut evaluasi seluruh kemungkinan koalisi fitur sehingga kompleksitasnya tumbuh secara eksponensial terhadap jumlah fitur. Oleh karena itu, pendekatan SHAP yang bersifat \emph{model-agnostic} umumnya hanya mengestimasi nilai Shapley melalui penarikan sampel, sehingga hasilnya bersifat aproksimatif dan tetap berbeban komputasi tinggi \autocite{Lundberg2017Unified}. Untuk mengatasi keterbatasan tersebut pada model berbasis pohon, \textcite{Lundberg2020TreeSHAP} mengembangkan TreeSHAP, yaitu algoritma yang dirancang khusus (\emph{model-specific}) untuk model berbasis pohon. Dengan memanfaatkan struktur pohon keputusan, TreeSHAP mampu menghitung nilai Shapley secara eksak (\emph{exact Shapley values}) dalam waktu polinomial, yaitu $O(TLD^2)$ dengan $T$ menyatakan jumlah pohon, $L$ jumlah daun maksimum pada suatu pohon, dan $D$ kedalaman maksimum pohon, alih-alih waktu eksponensial pada perhitungan langsung. Selain menghasilkan atribusi lokal yang eksak, \textcite{Lundberg2020TreeSHAP} juga menunjukkan bahwa nilai TreeSHAP dapat diagregasikan untuk memperoleh pemahaman global mengenai perilaku model, sehingga satu metode yang sama dapat menjelaskan prediksi baik pada tingkat individu maupun populasi.

Model yang digunakan pada penelitiann ini mayoritas merupakan model \emph{tree ensemble}. Berdasarkan landasan teori tersebut, metode interpretasi yang dipilih adalah TreeSHAP karena secara khusus dikembangkan untuk model berbasis pohon dan mampu menghitung nilai Shapley secara eksak sekaligus efisien \autocite{Lundberg2020TreeSHAP}. Pemilihan TreeSHAP memastikan bahwa atribusi fitur yang dihasilkan merupakan kontribusi yang tepat sesuai perilaku model, bukan hasil aproksimasi, sebuah syarat penting bagi sistem pendukung keputusan medis yang menuntut penjelasan yang akurat dan dapat dipertanggungjawabkan. Landasan teori pemilihan TreeSHAP ini menjadi acuan pada tahap perancangan dan implementasi modul interpretasi yang dibahas pada Bab \ref{chap:desain-konsep-solusi} dan Bab \ref{chap:hasil-pembahasan}.

% --- II.4 Large Language Model (LLM) ---
\section{Large Language Model (LLM)}
Large Language Model (LLM) merupakan evolusi mutakhir dalam domain Artificial Intelligence, khususnya dalam pemrosesan bahasa alami. LLM adalah model probabilitas yang dirancang untuk memahami, memproses, dan menghasilkan teks dengan tingkat koherensi yang menyerupai kemampuan manusia \autocite{Naveed2023Comprehensive}.

Arsitektur dasar dari teknologi ini mengandalkan teknik Deep Learning dan model statistik canggih untuk menganalisis korpus data teks dalam skala masif. Melalui pelatihan pada dataset berskala besar, model ini mampu mengidentifikasi pola linguistik yang kompleks, hubungan semantik antar kata, serta nuansa kontekstual yang terkandung dalam bahasa \autocite{Kasneci2023ChatGPT}.

Dalam ranah \emph{Generative AI}, kemampuan LLM tidak terbatas pada pemahaman teks semata, melainkan juga mencakup kemampuan untuk menciptakan konten orisinal yang mengadopsi gaya dan karakteristik data pelatihannya. Salah satu implementasi paling prominen dari teknologi ini adalah ChatGPT, sebuah model generatif yang dioptimalkan untuk dialog interaktif \autocite{OpenAI2023GPT4}.

Model-model ini dilatih menggunakan pendekatan \emph{Unsupervised Learning} atau \emph{Self-Supervised Learning}, yang memungkinkan mereka mempelajari struktur bahasa tanpa memerlukan pelabelan data manual secara ekstensif \autocite{Brown2020Language}. Fundamental teknologi di balik LLM adalah arsitektur Transformer, yang memperkenalkan mekanisme \emph{attention} untuk menangani ketergantungan jarak jauh dalam urutan data. Hal ini memungkinkan model untuk mempertahankan konteks pembicaraan atau dokumen yang panjang, menjadikannya sangat efektif untuk aplikasi seperti penerjemahan otomatis, \emph{chatbot} cerdas, dan pembuatan ringkasan teks otomatis \autocite{Vaswani2017Attention, Naveed2023Comprehensive}.


\subsection{Keterbatasan LLM}
Penerapan LLM di bidang medis menghadapi tantangan serius terkait fenomena halusinasi, yaitu kondisi saat model menghasilkan informasi yang terdengar meyakinkan namun faktanya salah atau dibuat-buat \autocite{Ji2023Survey}. Selain itu, LLM memiliki batasan pengetahuan statis atau \emph{knowledge cutoff} yang membuatnya tidak mengetahui informasi atau pedoman medis terbaru yang rilis setelah tanggal pelatihan model.

\subsection{Mengatasi Keterbatasan LLM}
Berbagai strategi dikembangkan untuk memitigasi keterbatasan tersebut. Salah satu pendekatan adalah \emph{fine-tuning} yang melatih ulang model dengan dataset spesifik, namun metode ini membutuhkan sumber daya komputasi yang besar dan data berlabel berkualitas tinggi. Pendekatan alternatif yang lebih efisien adalah \emph{In-Context Learning} melalui \emph{prompt engineering}, yaitu menyertakan informasi relevan secara langsung di dalam instruksi masukan tanpa mengubah bobot model. Kemampuan model bahasa berskala besar untuk mempelajari pola tugas hanya dari contoh dan konteks yang disertakan dalam \emph{prompt}, tanpa pembaruan parameter, pertama kali ditunjukkan secara empiris pada model GPT-3 \autocite{Brown2020Language}. Karena tidak memerlukan pelatihan ulang, pendekatan ini menghindari kebutuhan akan sumber daya komputasi besar maupun dataset berlabel spesifik-tugas yang menjadi prasyarat \emph{fine-tuning}, sehingga lebih ringan, murah, dan fleksibel untuk diterapkan pada domain spesifik \autocite{Liu2023Prompt}. Keunggulan efisiensi inilah yang menjadikan \emph{in-context learning} melalui \emph{prompt engineering} sebagai strategi yang dipilih dalam penelitian ini untuk mengarahkan keluaran LLM tanpa beban pelatihan ulang model.

% --- II.5 Prompt Engineering ---
\section{Prompt Engineering}
Prompt engineering merupakan disiplin perancangan instruksi masukan (\emph{prompt}) yang bertujuan mengarahkan perilaku Large Language Model (LLM) agar menghasilkan keluaran sesuai kebutuhan tugas tanpa mengubah parameter internal model \autocite{Brown2020Language}. Berbeda dengan \emph{fine-tuning} yang menuntut pelatihan ulang dan sumber daya komputasi besar, prompt engineering memanfaatkan kemampuan \emph{in-context learning} model, yaitu kemampuan LLM untuk menyesuaikan respons hanya berdasarkan konteks dan contoh yang disertakan di dalam prompt \autocite{Brown2020Language}. Pendekatan ini menjadikan prompt engineering sebagai metode yang ringan, murah, dan fleksibel untuk mengendalikan LLM pada aplikasi spesifik domain tanpa modifikasi bobot model.

\subsection{Teknik Perancangan Prompt}
Terdapat beberapa teknik dasar yang umum digunakan dalam perancangan prompt \autocite{Naveed2023Comprehensive}:
\begin{enumerate}
    \item \textbf{Zero-shot prompting}\\
    Model diberikan instruksi tugas secara langsung tanpa menyertakan contoh. Teknik ini mengandalkan pengetahuan yang telah dipelajari model selama pralatih dan sesuai untuk tugas yang relatif umum.

    \item \textbf{Few-shot prompting}\\
    Prompt dilengkapi sejumlah kecil contoh pasangan masukan--keluaran sebagai demonstrasi pola yang diharapkan. Contoh tersebut membantu model menyelaraskan format dan gaya jawaban dengan ekspektasi tugas \autocite{Brown2020Language}.

    \item \textbf{Role prompting}\\
    Model diberi peran atau persona tertentu (misalnya "Anda adalah dokter yang empatik") untuk mengarahkan gaya bahasa, sudut pandang, dan tingkat formalitas keluaran agar sesuai dengan konteks komunikasi kepada pasien.

    \item \textbf{Constraint atau guardrail prompting}\\
    Prompt menyertakan aturan eksplisit yang membatasi ruang keluaran model, seperti larangan mengarang angka, kewajiban hanya menggunakan fakta yang disediakan, serta batasan panjang dan format \autocite{Hakim2025Guardrails}. \emph{Guardrail} dalam konteks ini dipahami sebagai serangkaian batasan yang diterapkan pada masukan, proses, maupun keluaran LLM agar model tetap patuh pada kriteria yang telah ditetapkan \autocite{Hakim2025Guardrails}. Teknik ini krusial pada domain berisiko tinggi seperti kesehatan untuk menekan keluaran yang tidak diinginkan. \textcite{Hakim2025Guardrails} menunjukkan penerapannya pada setelan medis kritis-keselamatan (\emph{pharmacovigilance}), yaitu dengan merancang serangkaian \emph{guardrail} yang melarang model mengarang istilah obat maupun kejadian tidak diinginkan dan mewajibkan pengungkapan ketidakpastian, sehingga keluaran naratif LLM terkendali dan risiko halusinasi yang membahayakan pasien dapat ditekan.
\end{enumerate}

\subsection{Prompt Engineering sebagai Mitigasi Halusinasi}
Pada domain medis, tantangan utama penerapan LLM adalah fenomena halusinasi, yaitu keluaran yang koheren secara bahasa namun keliru secara fakta \autocite{Ji2023Survey}. Salah satu strategi mitigasi yang efektif tanpa memerlukan basis pengetahuan eksternal adalah membatasi peran LLM melalui prompt yang \emph{grounded}, yakni mengunci keluaran model pada fakta terstruktur yang telah dihasilkan oleh komponen lain dari sistem \autocite{Ji2023Survey}. Dalam konteks penelitian ini, fakta terstruktur tersebut berupa hasil atribusi fitur dari analisis SHAP. Dengan menyisipkan nilai kontribusi tiap faktor risiko secara eksplisit di dalam prompt serta melarang model menambahkan klaim di luar fakta tersebut, peran LLM dipersempit menjadi penerjemah dan penyampai edukasi gaya hidup umum, bukan sumber pengetahuan medis yang otonom. Strategi ini memungkinkan sistem memperoleh manfaat naratif LLM sekaligus menjaga akuntabilitas faktual, tanpa kompleksitas dan beban komputasi mekanisme temu-kembali (\emph{retrieval}).


% \section{System Usability Scale (SUS)}

% System Usability Scale (SUS) merupakan instrumen evaluasi usability yang diperkenalkan oleh Brooke untuk mengukur persepsi pengguna terhadap kemudahan penggunaan suatu sistem \parencite{Brooke1996}. SUS dirancang sebagai metode evaluasi yang sederhana, cepat, dan dapat diterapkan pada berbagai jenis produk digital, termasuk perangkat lunak, aplikasi web, sistem informasi, maupun perangkat keras. Instrumen ini terdiri atas sepuluh pernyataan yang dinilai menggunakan skala Likert lima poin, mulai dari `Sangat Tidak Setuju'' hingga `Sangat Setuju'' \parencite{Brooke1996}.

% Dalam SUS, pernyataan bernomor ganjil merepresentasikan aspek positif terhadap usability, sedangkan pernyataan bernomor genap merepresentasikan aspek negatif. Skor SUS dihitung dengan melakukan normalisasi terhadap setiap respons, kemudian menjumlahkan seluruh skor dan mengalikannya dengan faktor 2,5 sehingga menghasilkan nilai akhir pada rentang 0 hingga 100 \parencite{Brooke1996}. Meskipun berada pada rentang 0 hingga 100, skor SUS tidak merepresentasikan persentase tingkat usability, melainkan indeks yang digunakan untuk membandingkan tingkat kegunaan suatu sistem.

% Berbagai penelitian menunjukkan bahwa SUS memiliki tingkat reliabilitas yang tinggi dan mampu memberikan hasil evaluasi yang konsisten meskipun hanya terdiri atas sepuluh item pertanyaan. Selain itu, SUS telah digunakan secara luas pada berbagai domain aplikasi sehingga memungkinkan hasil evaluasi suatu sistem dibandingkan dengan hasil penelitian lainnya \parencite{Bangor2008}.

% Bangor \textit{et al.} melakukan evaluasi empiris terhadap ribuan respons SUS dan menunjukkan bahwa skor SUS dapat diinterpretasikan menggunakan kategori tertentu untuk mempermudah analisis usability. Secara umum, skor di atas rata-rata 68 menunjukkan usability yang baik, sedangkan skor di bawah nilai tersebut mengindikasikan perlunya perbaikan pada sistem \parencite{Bangor2008}. Interpretasi ini menjadikan SUS tidak hanya berfungsi sebagai instrumen pengukuran, tetapi juga sebagai dasar dalam mengevaluasi kualitas pengalaman pengguna secara kuantitatif.

% Untuk pengguna berbahasa Indonesia, Sharfina dan Santoso melakukan adaptasi SUS ke dalam Bahasa Indonesia melalui proses penerjemahan dan pengujian reliabilitas. Hasil penelitian menunjukkan bahwa versi Bahasa Indonesia dari SUS memiliki karakteristik psikometrik yang baik dan dapat digunakan sebagai instrumen evaluasi usability pada populasi pengguna Indonesia \parencite{Sharfina2016}. Oleh karena itu, penelitian ini menggunakan versi Bahasa Indonesia dari SUS sebagai instrumen untuk mengevaluasi tingkat kemudahan penggunaan sistem informasi berbasis Machine Learning dan Explainable Artificial Intelligence yang dikembangkan.

% \begin{table}[ht]
% \centering
% \caption{Versi Bahasa Indonesia dari System Usability Scale (SUS)}
% \label{tab:sus-indonesia}
% \begin{tabular}{cp{11cm}}
% \hline
% No. & Pernyataan \\ \hline
% 1 & Saya berpikir akan menggunakan sistem ini lagi. \\
% 2 & Saya merasa sistem ini rumit untuk digunakan. \\
% 3 & Saya merasa sistem ini mudah untuk digunakan. \\
% 4 & Saya membutuhkan bantuan dari orang lain atau teknisi dalam menggunakan sistem ini. \\
% 5 & Saya merasa fitur-fitur sistem ini berjalan dengan semestinya. \\
% 6 & Saya merasa ada banyak hal yang tidak konsisten (tidak serasi) pada sistem ini. \\
% 7 & Saya merasa orang lain akan memahami cara menggunakan sistem ini dengan cepat. \\
% 8 & Saya merasa sistem ini membingungkan. \\
% 9 & Saya merasa tidak ada hambatan dalam menggunakan sistem ini. \\
% 10 & Saya perlu membiasakan diri terlebih dahulu sebelum menggunakan sistem ini. \\
% \hline
% \end{tabular}

% \vspace{1mm}
% \footnotesize
% Sumber: Adaptasi dari \textcite{Sharfina2016}.
% \end{table}
% --- II.6 Tinjauan Penelitian Sebelumnya ---
\section{Tinjauan Penelitian Sebelumnya}
Penelitian mengenai prediksi risiko hipertensi telah berkembang pesat dalam beberapa tahun terakhir, bergeser dari metode statistik konvensional menuju pendekatan Machine Learning (ML) yang lebih mampu menangani data tabular dengan variabel yang beragam. Fokus utama berbagai studi adalah meningkatkan akurasi deteksi dini agar intervensi dapat dilakukan sebelum komplikasi kardiovaskular terjadi \autocite{Silva2022MLHypertension}. Metode berbasis \emph{ensemble} seperti Random Forest dan varian \emph{gradient boosting} secara konsisten dilaporkan mengungguli model tunggal pada dataset kesehatan yang tidak seimbang. Dalam konteks Indonesia, sejumlah studi bahkan telah memanfaatkan data berskala besar untuk membangun model skrining hipertensi berbasis faktor risiko non-laboratorium \autocite{Estiko2025, Septian2025Prediction}.

Meskipun akurasinya tinggi, model ML yang kompleks umumnya bersifat \emph{black box} sehingga alasan di balik setiap prediksinya sulit dipahami. Untuk mengatasi keterbatasan ini, gelombang penelitian berikutnya mengintegrasikan metode Explainable Artificial Intelligence (XAI) seperti SHAP dan LIME guna mengungkap kontribusi tiap fitur terhadap hasil prediksi \autocite{Adadi2018}. Melalui atribusi fitur tersebut, sistem tidak hanya memberikan label risiko, tetapi juga menunjukkan variabel mana, misalnya usia, lingkar pinggang, atau riwayat keluarga, yang paling memengaruhi keputusan model. Namun, keluaran XAI umumnya masih berupa nilai numerik atau grafik yang sulit dicerna oleh masyarakat awam.

Di sisi lain, kemunculan Large Language Model (LLM) membuka peluang baru untuk menjembatani kesenjangan komunikasi tersebut. LLM memiliki potensi besar dalam merangkum informasi medis yang kompleks menjadi penjelasan bahasa alami yang personal dan mudah dipahami \autocite{Thirunavukarasu2023Large}. Meskipun demikian, penerapan LLM di ranah medis masih dibayangi risiko halusinasi, yaitu keluaran yang terdengar meyakinkan namun keliru secara fakta, sehingga penggunaannya perlu dikendalikan secara ketat \autocite{Clusmann2023Future}. Belakangan mulai bermunculan studi yang memadukan atribusi XAI dengan LLM untuk menghasilkan narasi penjelasan, misalnya kerangka \emph{prompt} yang menerjemahkan nilai SHAP menjadi penjelasan naratif \autocite{Lee2025Prompt}. Meskipun demikian, upaya integrasi semacam ini masih terbatas pada domain klinis lain seperti perawatan intensif dan penyakit kardiovaskular, serta umumnya ditujukan bagi tenaga medis, bukan masyarakat awam.

Berdasarkan tinjauan tersebut, terdapat celah penelitian yang dapat diisi. Pada ranah hipertensi, mayoritas studi berhenti pada tahap prediksi ML dan interpretasi XAI yang keluarannya masih bersifat teknis, sedangkan integrasi lanjutan dengan LLM untuk menerjemahkan hasil tersebut menjadi narasi yang dapat dipahami masyarakat awam belum ditemukan, khususnya pada konteks populasi Indonesia dengan fitur non-invasif. Meskipun pendekatan serupa mulai berkembang di domain klinis lain, penelitian ini menerapkannya secara spesifik untuk estimasi risiko hipertensi dengan memadukan model prediksi ML, interpretasi SHAP, dan LLM yang dikendalikan melalui \emph{prompt engineering}, disertai pengendalian halusinasi tanpa basis pengetahuan eksternal, guna menghasilkan estimasi yang akurat, transparan, sekaligus komunikatif bagi pengguna awam. Untuk memberikan gambaran posisi penelitian ini dibandingkan studi terdahulu, Tabel \ref{tbl:studi_sebelumnya} merangkum lima penelitian relevan beserta metode, hasil, dan keterbatasannya.

% Pengaturan agar tabel mentok kiri-kanan (Full Width)
\setlength\LTleft{0pt} 
\setlength\LTright{0pt}

% Pengaturan agar tabel naik ke atas (Mentok Header / Mengurangi spasi atas)
\setlength{\LTpre}{0pt} 
\setlength{\LTpost}{0pt}

\begin{longtable}{@{\extracolsep{\fill}} p{0.5cm} p{3.5cm} p{3.5cm} p{3.5cm} p{3.5cm}}
\caption{Studi Sebelumnya} \label{tbl:studi_sebelumnya} \\
\toprule
\textbf{No} & \textbf{Judul \& Penulis} & \textbf{Metode} & \textbf{Hasil} & \textbf{Keterbatasan} \\
\midrule
\endfirsthead


\toprule
\textbf{No} & \textbf{Judul \& Penulis} & \textbf{Metode} & \textbf{Hasil} & \textbf{Keterbatasan} \\
\midrule
\endhead

\midrule
\multicolumn{5}{r}{\emph{Bersambung ke halaman berikutnya}} \\
\endfoot

\bottomrule
\endlastfoot

1 & Screening Hypertension Using Non-Laboratory Risk Factors With Machine Learning: A Retrospective Cross-Sectional Study in Indonesia \autocite{Estiko2025} &
Membandingkan Random Forest, Gradient Boosting, dan XGBoost dengan Logistic Regression sebagai \emph{baseline} untuk menyaring hipertensi menggunakan sebelas faktor risiko non-laboratorium pada ratusan ribu peserta Posbindu di Indonesia, dengan validasi internal (\emph{10-fold cross-validation}) dan eksternal. &
Model berbasis \emph{boosting} (XGBoost) berkinerja terbaik dan mengungguli \emph{baseline} Logistic Regression pada validasi internal maupun eksternal. Riwayat keluarga hipertensi, usia, lingkar pinggang, dan IMT menjadi faktor prediktor paling berpengaruh. &
Bersifat \emph{cross-sectional} sehingga tidak dapat menyimpulkan kausalitas atau insidensi, hanya mengandalkan faktor non-laboratorium, dan hasil prediksinya belum diterjemahkan menjadi penjelasan naratif bagi pengguna awam. \\
\hline
2 & Prediction of Personalised Hypertension Using Machine Learning in Indonesian Population \autocite{Septian2025Prediction} &
Mengomparasi lima algoritma, yaitu XGBoost, LightGBM, CatBoost, Logistic Regression, dan Random Forest, pada 9,58 juta data SATUSEHAT dengan analisis kontribusi fitur menggunakan SHAP (\emph{SHapley Additive exPlanations}). &
Model LightGBM mencapai performa terbaik dengan AUC 0,78 tanpa riwayat pasien dan 0,85 dengan riwayat pasien. Analisis SHAP menunjukkan usia, riwayat keluarga, berat badan, dan lingkar pinggang sebagai prediktor dominan. &
Studi bersifat \emph{cross-sectional} sehingga tidak dapat menyimpulkan kausalitas, terbatas pada seting non-rumah sakit tanpa variabel klinis mendalam, serta hasil interpretasinya belum disajikan sebagai narasi untuk awam. \\
\hline
3 & XGBoost and SHAP-Based Analysis of Risk Factors for Hypertension Classification in Korean Postmenopausal Women \autocite{Li2025XGBoost} &
Menggunakan XGBoost, Support Vector Machine (SVM), dan Artificial Neural Network (ANN) untuk klasifikasi hipertensi, dengan interpretabilitas model melalui SHAP guna mengidentifikasi faktor risiko paling berpengaruh. &
Model XGBoost memberikan performa terbaik dengan AUC 92,12\%, akurasi 84,73\%, dan MCC 0,71. SHAP mengidentifikasi usia dan lingkar pinggang sebagai faktor paling dominan. &
Terbatas pada populasi wanita pascamenopause di Korea Selatan sehingga sulit digeneralisasi, serta faktor gaya hidup seperti konsumsi garam, aktivitas fisik, dan tingkat stres tidak disertakan karena keterbatasan data. \\
\hline
4 & Explainable AI in Action: A Comparative Analysis of Hypertension Risk Factors Using SHAP and LIME \autocite{Donmez2025} &
Menerapkan model XGBoost untuk memprediksi hipertensi pada 623 pasien dengan 13 biomarker, lalu membandingkan dua metode XAI, yaitu SHAP dan LIME, untuk menjelaskan kontribusi fitur terhadap prediksi. &
Model XGBoost mencapai akurasi 99,4\%, presisi 100\%, recall 97,30\%, dan F1-score 98,6\%. SHAP dan LIME berhasil mengungkap faktor risiko utama di balik prediksi model. &
Dataset relatif kecil (623 pasien) dengan metrik nyaris sempurna yang berpotensi \emph{overfitting}, bergantung pada biomarker laboratorium sehingga tidak non-invasif, dan penjelasannya belum disajikan dalam bentuk narasi bahasa alami. \\
\hline
5 & Large Language Models in Medicine \autocite{Thirunavukarasu2023Large} &
Tinjauan mengenai penerapan Large Language Model (LLM) berbasis arsitektur Transformer untuk beragam tugas klinis, termasuk komunikasi dan edukasi kepada pasien. &
LLM berpotensi menyintesis informasi medis kompleks menjadi penjelasan yang mudah dipahami oleh pasien maupun tenaga medis. &
Rentan terhadap halusinasi fakta dan bias data pelatihan, serta belum diintegrasikan dengan model prediksi risiko numerik seperti keluaran ML dan XAI. \\
\end{longtable}